\documentclass[12pt, a4paper]{exam}
% Corrigé rédigé pour le site Ressources NSI et SNT (CC BY-SA 4.0).
% Compilation : pdflatex (deux passes).
\usepackage[T1]{fontenc}
\usepackage[utf8]{inputenc}
\usepackage{lmodern}
\usepackage[french]{babel}
\usepackage[a4paper,top=2.3cm,bottom=2.8cm,left=2cm,right=2cm,headsep=0.6cm,footskip=1.2cm]{geometry}
\usepackage{xcolor}
\usepackage{listings}
\usepackage{graphicx}
\usepackage{array}
\usepackage{amsmath}
\usepackage{amssymb}
\usepackage{enumitem}
\usepackage{needspace}
\usepackage{tikz}
\usetikzlibrary{positioning,arrows.meta,calc}
\usepackage[most]{tcolorbox}

% Coloration du code : même style que les sujets
\definecolor{codegreen}{rgb}{0,0.6,0}
\definecolor{codegray}{rgb}{0.5,0.5,0.5}
\definecolor{codepurple}{rgb}{0.58,0,0.82}
\definecolor{backcolour}{rgb}{0.95,0.95,0.92}
\lstdefinestyle{mystyle}{
    backgroundcolor=\color{backcolour},
    commentstyle=\color{codegreen}\itshape,
    keywordstyle=\color{magenta}\bfseries,
    numberstyle=\tiny\color{codegray},
    stringstyle=\color{codepurple},
    basicstyle=\footnotesize\ttfamily,
    breakatwhitespace=false,
    breaklines=true,
    postbreak=\mbox{\textcolor{codegray}{$\hookrightarrow$}\space},
    captionpos=b,
    keepspaces=true,
    columns=fullflexible,
    numbers=left,
    numbersep=5pt,
    xleftmargin=1.5em,
    showspaces=false,
    showstringspaces=false,
    showtabs=false,
    tabsize=4,
    upquote=true,
    extendedchars=true,
    literate={à}{{\`a}}1 {â}{{\^a}}1 {é}{{\'e}}1 {è}{{\`e}}1 {ê}{{\^e}}1 {ë}{{\"e}}1
             {î}{{\^i}}1 {ï}{{\"i}}1 {ô}{{\^o}}1 {ù}{{\`u}}1 {û}{{\^u}}1 {ç}{{\c c}}1
             {É}{{\'E}}1 {À}{{\`A}}1 {œ}{{\oe}}1 {…}{{\ldots}}1
}
\lstset{style=mystyle, language=Python}
\newcommand{\code}[1]{\texttt{#1}}

% Encadrés du corrigé (lisibles en noir et blanc)
\tcbset{corrbase/.style={enhanced, breakable, sharp corners, left=2mm, right=2mm, top=1mm, bottom=1mm,
  fonttitle=\bfseries\small, coltitle=black, before skip=2mm, after skip=3mm}}
\newtcolorbox{reponse}[1][Réponse]{corrbase, unbreakable, colback=gray!6, colframe=black, colbacktitle=gray!22, boxrule=0.8pt, title={#1}}
\newtcolorbox{erreur}[1][Erreur fréquente]{corrbase, colback=white, colframe=black, boxrule=0.4pt, leftrule=3pt, title={#1}, colbacktitle=white, attach title to upper={\ :\ }}
\newtcolorbox{retenir}[1][À retenir]{corrbase, colback=white, colframe=black, boxrule=0.4pt, borderline={0.4pt}{1.2pt}{black}, title={#1}, colbacktitle=white, attach title to upper={\ :\ }}

% En-tête d'une question : numéro, points, rappel de la question
\newcommand{\Q}[3]{\par\needspace{5\baselineskip}\bigskip\noindent{\large\bfseries Question #1}\ifx&#2&\else\ \textnormal{(#2)}\fi\par\nopagebreak\smallskip\noindent\textit{#3}\par\nopagebreak\smallskip}
\newcommand{\Qp}[2]{\par\medskip\noindent\textbf{#1}\ \textit{#2}\par\nopagebreak\smallskip}
\newcommand{\methode}{\par\noindent\textbf{Méthode.}\ }
\newcommand{\application}{\par\noindent\textbf{Application.}\ }

% Arbres
\tikzset{arbre/.style={every node/.style={draw,circle,minimum size=7mm,inner sep=1pt,font=\small},level distance=11mm, thick},
  marque/.style={fill=gray!35, line width=1.6pt},
  titre/.style={draw=none, rectangle, inner sep=2pt}}

\raggedbottom
\setlength{\parindent}{0pt}
\setlength{\parskip}{3pt}
\setlist{itemsep=1pt, topsep=2pt}
\newtcolorbox{remarque}[1][Remarque sur l'énoncé]{corrbase, colback=white, frame hidden, borderline={0.6pt}{0pt}{black, dashed}, title={#1}, colbacktitle=white, attach title to upper={\ :\ }}

\title{Corrigé --- CC Term NSI - Objet}
\author{Babash - NSI Tani Malandi et Johanne Bakalara}
\date{Octobre 2024 \textperiodcentered{} Durée du contrôle : 1 h 30 \textperiodcentered{} 23 points}

% Une maison et une antenne sur l'axe de la rue
\newcommand{\maison}[1]{\begin{scope}[shift={(#1,0.12)}]\draw (-0.22,0) rectangle (0.22,0.42);\draw (-0.22,0.42) -- (0,0.68) -- (0.22,0.42);\end{scope}}
\newcommand{\antenne}[3]{% position, rayon, hauteur de la portée
\draw[thick] (#1,-0.05) -- (#1-0.12,-0.75) -- (#1+0.12,-0.75) -- cycle;
\draw[|<->|, thick] (#1-#2,-1.0-#3) -- (#1+#2,-1.0-#3);
\node[font=\scriptsize, fill=white, inner sep=1pt] at (#1,-1.0-#3) {#1};}
\newcommand{\axe}{%
\draw[very thick] (-0.6,0) -- (13.2,0);
\foreach \i in {0,...,13} {\draw (\i,-0.08) -- (\i,0.08);}
\foreach \i in {0,5,10} {\node[above, font=\scriptsize] at (\i,0.75) {\i};}
\foreach \p in {0,2,3,4,5,7,9,10.5,11.5} {\maison{\p}}}

\begin{document}
\pagestyle{foot}
\footer{}{\footnotesize Babash - NSI Tani Malandi et Johanne Bakalara \textperiodcentered{} CC BY-SA 4.0 \textperiodcentered{} page \thepage}{}
\maketitle
\thispagestyle{foot}
\begin{center}\footnotesize\itshape Exercice des antennes relais d'après l'exercice 1 du sujet 24-NSIJ1JA1 (baccalauréat 2024, jour 1).\end{center}

\begin{retenir}[Comment lire ce corrigé]
Pour chaque question : la question rappelée en italique, la réponse attendue dans un cadre gris, puis l'explication. Le code de chaque réponse a été exécuté. Les cadres « Erreur fréquente » montrent les pièges classiques. Le barème du sujet totalise 23 points.
\end{retenir}

\begin{retenir}[Le vocabulaire de la programmation objet]
\begin{itemize}
\item Une \textbf{classe} est un modèle : elle décrit ce que contient et ce que sait faire un type d'objet (par exemple \code{Eleve}).
\item Un \textbf{objet}, ou \textbf{instance}, est un exemplaire fabriqué à partir de la classe (par exemple l'élève Alex).
\item Un \textbf{attribut} est une variable rangée dans l'objet : on l'écrit \code{self.nom\_attribut}.
\item Une \textbf{méthode} est une fonction définie dans la classe ; elle reçoit l'objet lui-même dans le paramètre \code{self}.
\item Le \textbf{constructeur} \code{\_\_init\_\_} est la méthode appelée automatiquement à la création d'un objet ; il crée les attributs.
\end{itemize}
\end{retenir}

\section*{1. Cours}

Le code de la classe \code{Eleve}, avec ses numéros de ligne :
\begin{lstlisting}
class Eleve:
    def __init__(self, nom, prenom, notes):
        self.name = nom
        self.first_name = prenom
        for n in notes.values():
            assert 0 <= n <= 20
        self.grades = notes

    def moyenne_sans_coeff(self):
        somme = 0
        for n in self.grades.values():
            somme += n
        return somme/len(self.grades)

    def moyenne_coeff(self, coeff):
        somme = 0
        somme_coeff = 0
        for m, n in self.grades.items():
            somme += coeff[m] * n
            somme_coeff += coeff[m]
        return somme/somme_coeff
\end{lstlisting}

\Q{1}{4 points}{(a) Attributs de la classe et leur type. (b) Méthodes de la classe. (c) Que fait la ligne 6 ? (d) Que signifie POO ?}
\begin{reponse}
(a) Trois attributs : \code{name} (chaîne de caractères, \code{str}), \code{first\_name} (chaîne de caractères, \code{str}) et \code{grades} (dictionnaire, \code{dict}, qui associe à chaque matière une note). \\
(b) Trois méthodes : \code{\_\_init\_\_} (le constructeur), \code{moyenne\_sans\_coeff} et \code{moyenne\_coeff}. \\
(c) La ligne 6 vérifie que chaque note \code{n} est comprise entre 0 et 20. Si une note est en dehors, le programme s'arrête avec une erreur \code{AssertionError}, et l'objet n'est pas créé. \\
(d) POO signifie \textbf{programmation orientée objet}.
\end{reponse}
\begin{itemize}
\item (a) Les attributs se repèrent à \code{self.} suivi d'un nom, dans le constructeur (lignes 3, 4 et 7). Attention : les attributs s'appellent \code{name}, \code{first\_name} et \code{grades}, et non \code{nom}, \code{prenom} et \code{notes}, qui sont les \emph{paramètres} du constructeur. Le type de \code{grades} se déduit de la ligne 5 : \code{notes.values()} n'existe que pour un dictionnaire.
\item (b) Les méthodes sont les \code{def} écrits \emph{dans} la classe (indentés sous \code{class Eleve:}).
\item (c) \code{assert condition} se lit « affirmer que la condition est vraie ». Si elle est vraie, rien ne se passe. Si elle est fausse, Python lève une erreur. Exemple : \code{Eleve("X", "Y", \{"NSI": 21\})} provoque \code{AssertionError}. La boucle de la ligne 5 fait ce test pour chaque note.
\item (d) En POO, on regroupe dans un même objet des données (les attributs) et les actions sur ces données (les méthodes).
\end{itemize}
\begin{erreur}
Citer \code{somme}, \code{n} ou \code{coeff} comme attributs. Ce sont des variables \emph{locales} ou des paramètres des méthodes : elles n'ont pas de \code{self.} devant et disparaissent à la fin de la méthode.
\end{erreur}

\section*{2. Travail sur les classes}

\Q{2 (a)}{}{Créer l'objet correspondant à Alex Martin (18 en NSI, 15 en anglais, 12 en arabe) dans la variable \code{alex}.}
\begin{reponse}
\begin{lstlisting}
alex = Eleve("Martin", "Alex", {"NSI": 18, "Anglais": 15, "Arabe": 12})
\end{lstlisting}
\end{reponse}
On appelle la classe comme une fonction, avec les paramètres du constructeur \textbf{dans l'ordre} : \code{nom}, puis \code{prenom}, puis \code{notes}. Le paramètre \code{self} ne s'écrit pas : Python le fournit tout seul. Les notes forment un dictionnaire : les clés sont les matières (des chaînes), les valeurs les notes.

\textbf{Barème.} La question 2 vaut 4 points pour ses quatre parties (a) à (d).

\Q{2 (b)}{}{Calculer la moyenne sans coefficient d'Alex.}
\begin{reponse}
\code{alex.moyenne\_sans\_coeff()} \quad (le résultat est \textbf{15.0})
\end{reponse}
On appelle une méthode avec un point : \code{objet.methode()}. Les parenthèses sont obligatoires, même vides. Calcul : $(18 + 15 + 12) / 3 = 45 / 3 = 15$.

\Q{2 (c)}{}{Calculer la moyenne avec coefficients : 6 pour la NSI, 3 pour l'anglais, 3 pour l'arabe.}
\begin{reponse}
\code{alex.moyenne\_coeff(\{"NSI": 6, "Anglais": 3, "Arabe": 3\})} \\
Le résultat est \textbf{15.75}.
\end{reponse}
La méthode utilise \code{coeff[m]}, où \code{m} est une matière de \code{grades} : les coefficients doivent donc être dans un dictionnaire \textbf{avec les mêmes clés} que les notes.
\textbf{Trace de la boucle} (lignes 18 à 20) :
\begin{center}
\renewcommand{\arraystretch}{1.15}
\begin{tabular}{|l|c|c|c|c|}
\hline
\code{m} & \code{n} & \code{coeff[m]} & \code{somme} & \code{somme\_coeff} \\ \hline
\code{"NSI"} & 18 & 6 & $0 + 6 \times 18 = 108$ & 6 \\
\code{"Anglais"} & 15 & 3 & $108 + 45 = 153$ & 9 \\
\code{"Arabe"} & 12 & 3 & $153 + 36 = 189$ & 12 \\ \hline
\end{tabular}
\end{center}
La méthode renvoie $189 / 12 = 15{,}75$.
\begin{erreur}
Écrire les clés avec une autre orthographe (\code{"nsi"}, \code{"anglais"}) : \code{coeff[m]} ne trouve pas la clé, et Python s'arrête sur une \code{KeyError} (« erreur de clé »). Les clés de deux dictionnaires doivent être écrites exactement de la même façon.
\end{erreur}

\Q{2 (d)}{}{Modifier la note de NSI d'Alex pour lui ajouter un point.}
\begin{reponse}
\code{alex.grades["NSI"] = alex.grades["NSI"] + 1}
\end{reponse}
\code{alex.grades} est le dictionnaire des notes ; \code{alex.grades["NSI"]} est la note de NSI (18). On la remplace par $18 + 1 = 19$. Écriture équivalente : \code{alex.grades["NSI"] += 1}.

\Q{3}{2 points}{Écrire une méthode \code{meilleure\_note} qui renvoie la meilleure note d'un·e élève et la matière où elle a été obtenue.}
\begin{reponse}
\begin{lstlisting}
    def meilleure_note(self):
        meilleure_matiere = None
        meilleure = -1
        for matiere, note in self.grades.items():
            if note > meilleure:
                meilleure = note
                meilleure_matiere = matiere
        return meilleure, meilleure_matiere
\end{lstlisting}
\end{reponse}
\methode C'est la recherche d'un maximum, vue en première :
\begin{enumerate}
\item On part d'une « meilleure note » plus petite que toutes les notes possibles : $-1$ (les notes sont entre 0 et 20).
\item On parcourt les couples (matière, note) avec \code{.items()}.
\item Si une note dépasse la meilleure trouvée, on retient cette note \textbf{et} sa matière.
\item À la fin, on renvoie les deux valeurs, dans un tuple.
\end{enumerate}
La méthode s'écrit \textbf{dans} la classe : elle est indentée sous \code{class Eleve:} et son premier paramètre est \code{self}.

\textbf{Trace} avec \code{alex} (notes 19 en NSI après la question 2 (d), 15, 12) : après « NSI », \code{meilleure} vaut 19 ; 15 et 12 ne sont pas plus grands. L'appel \code{alex.meilleure\_note()} renvoie \code{(19, 'NSI')}.

\textbf{Variante}, avec la fonction \code{max} et une clé (plus courte, mais hors programme) :
\begin{lstlisting}
    def meilleure_note(self):
        matiere = max(self.grades, key=self.grades.get)
        return self.grades[matiere], matiere
\end{lstlisting}
\begin{erreur}
Oublier \code{self} : \code{def meilleure\_note():} provoque une \code{TypeError} à l'appel, car Python passe toujours l'objet en premier paramètre. Autre erreur : ne renvoyer que la note, alors que la matière est aussi demandée.
\end{erreur}

\section*{3. Exercice type bac : les antennes relais}

\Q{4}{2 points}{Créer les variables \code{m1} et \code{m2}, instances de \code{Maison} aux abscisses 1 et 3,5.}
\begin{reponse}
\begin{lstlisting}
m1 = Maison(1)
m2 = Maison(3.5)
\end{lstlisting}
\end{reponse}
Le constructeur de \code{Maison} attend un seul paramètre : \code{position}. Attention, en Python, la virgule décimale s'écrit avec un \textbf{point} : \code{3.5}. Écrire \code{Maison(3,5)} donnerait deux paramètres et provoquerait une \code{TypeError}.

\Q{5}{1 point}{Créer la variable \code{a} : l'antenne à la position 2,5, de rayon d'action 1.}
\begin{reponse}
\code{a = Antenne(2.5, 1)}
\end{reponse}
Le constructeur de \code{Antenne} attend la \code{position} puis le \code{rayon}, dans cet ordre. \code{Antenne(1, 2.5)} créerait une antenne en 1 de rayon 2,5 : une autre antenne.

\Q{6}{1 point}{Représenter la situation créée par \code{creation\_rue([0, 2, 3, 4, 5, 7, 9, 10.5, 11.5])}.}
\begin{reponse}
\begin{center}
\begin{tikzpicture}[x=1.0cm]
\axe
\end{tikzpicture}
\end{center}
Neuf maisons, placées aux abscisses 0, 2, 3, 4, 5, 7, 9, 10,5 et 11,5.
\end{reponse}
Chaque nombre de la liste est l'abscisse d'une maison : on dessine une maison au-dessus de chaque graduation correspondante. Pour 10,5 et 11,5, la maison est au milieu entre deux graduations.

\Q{7}{2 points}{Compléter la fonction \code{creation\_rue}.}
\begin{reponse}
\begin{lstlisting}
def creation_rue(pos):
    maisons = []
    for p in pos:
        m = Maison(p)
        maisons.append(m)
    return maisons
\end{lstlisting}
\end{reponse}
Ligne 2 : on part d'une liste vide. Lignes 3 à 5 : pour chaque position \code{p}, on crée une maison et on l'ajoute à la fin de la liste. Ligne 6 : on renvoie la liste des objets \code{Maison}.
\begin{erreur}
Écrire \code{maisons.append(p)} : on ajouterait le \emph{nombre} \code{p}, pas l'objet \code{Maison}. La suite (\code{m.get\_pos\_maison()}) provoquerait alors une \code{AttributeError}, car un nombre n'a pas cette méthode.
\end{erreur}

\Q{8}{2 points}{Compléter la méthode \code{couvre} de la classe \code{Antenne}.}
\begin{reponse}
\begin{lstlisting}
def couvre(self, maison):
    return abs(maison.get_pos_maison() - self._position) <= self._rayon
\end{lstlisting}
\end{reponse}
\methode
\begin{enumerate}
\item La \textbf{distance} entre deux points d'un axe est la valeur absolue de la différence de leurs abscisses : $|x_{\text{maison}} - x_{\text{antenne}}|$.
\item L'antenne couvre la maison si cette distance est \textbf{inférieure ou égale} au rayon.
\item La comparaison donne déjà \code{True} ou \code{False} : on la renvoie directement.
\end{enumerate}
\application avec \code{a = Antenne(2.5, 1)} : pour \code{m1}, $|1 - 2{,}5| = 1{,}5 > 1$, donc \code{a.couvre(m1)} renvoie \code{False} ; pour \code{m2}, $|3{,}5 - 2{,}5| = 1 \leq 1$, donc \code{a.couvre(m2)} renvoie \code{True}. Ce sont bien les résultats annoncés.

\textbf{Variante}, avec les accesseurs : \code{abs(maison.get\_pos\_maison() - self.get\_pos\_antenne()) <= self.get\_rayon()}.
\begin{erreur}
Écrire \code{maison.\_position} : cela marche en Python, mais l'attribut commence par \code{\_}, ce qui signifie « privé » : depuis une autre classe, on passe par la méthode \code{get\_pos\_maison()}. Autre erreur : utiliser \code{<} au lieu de \code{<=}. L'énoncé dit « inférieure ou égale » : la maison \code{m2}, à distance exactement 1, ne serait plus couverte.
\end{erreur}

\Q{9}{2 points}{Que renvoie (affiche) la suite d'instructions avec \code{strategie\_1(maisons, 2)} ?}
\begin{reponse}
\code{[0, 3, 7, 10.5]}
\end{reponse}
\textbf{Le principe de} \code{strategie\_1} : on place une antenne sur la première maison. Puis, pour chaque maison suivante, si la \textbf{dernière} antenne posée ne la couvre pas, on pose une nouvelle antenne \textbf{sur} cette maison.

\textbf{Tableau de trace} (rayon 2 ; une antenne en $x$ couvre de $x - 2$ à $x + 2$) :
\begin{center}
\renewcommand{\arraystretch}{1.15}
\begin{tabular}{|c|c|l|l|}
\hline
\textbf{Maison} & \textbf{Dernière antenne} & \textbf{Couverte ?} & \textbf{Antennes} \\ \hline
0 & --- & première maison : antenne en 0 & \code{[0]} \\
2 & 0 & $|2 - 0| = 2 \leq 2$ : oui & \code{[0]} \\
3 & 0 & $|3 - 0| = 3 > 2$ : non, antenne en 3 & \code{[0, 3]} \\
4 & 3 & $|4 - 3| = 1$ : oui & \code{[0, 3]} \\
5 & 3 & $|5 - 3| = 2$ : oui & \code{[0, 3]} \\
7 & 3 & $|7 - 3| = 4 > 2$ : non, antenne en 7 & \code{[0, 3, 7]} \\
9 & 7 & $|9 - 7| = 2$ : oui & \code{[0, 3, 7]} \\
10,5 & 7 & $3{,}5 > 2$ : non, antenne en 10,5 & \code{[0, 3, 7, 10.5]} \\
11,5 & 10,5 & $1 \leq 2$ : oui & \code{[0, 3, 7, 10.5]} \\ \hline
\end{tabular}
\end{center}
La dernière ligne affiche la liste des positions des antennes, construite par compréhension : \code{[0, 3, 7, 10.5]}. Il faut \textbf{4 antennes}.

\Q{10}{1 point}{Placer les antennes selon la nouvelle stratégie (rayon 2) sur la rue [0, 2, 3, 4, 5, 7, 9, 10.5, 11.5].}
\begin{reponse}
Trois antennes, aux abscisses \textbf{2}, \textbf{7} et \textbf{12,5} :
\begin{center}
\begin{tikzpicture}[x=1.0cm]
\axe
\antenne{2}{2}{0}
\antenne{7}{2}{0}
\antenne{12.5}{2}{0}
\end{tikzpicture}
\end{center}
\end{reponse}
\textbf{La nouvelle règle} : dès qu'une maison n'est pas couverte, on place l'antenne le plus à droite possible tout en couvrant cette maison, c'est-à-dire à l'abscisse de la maison $+ \, r$.
\begin{enumerate}
\item La maison 0 n'est pas couverte : antenne en $0 + 2 = 2$. Elle couvre de 0 à 4 : maisons 0, 2, 3 et 4.
\item La maison 5 n'est pas couverte ($|5 - 2| = 3 > 2$) : antenne en $5 + 2 = 7$. Elle couvre de 5 à 9 : maisons 5, 7 et 9.
\item La maison 10,5 n'est pas couverte : antenne en $10{,}5 + 2 = 12{,}5$. Elle couvre de 10,5 à 14{,}5 : maisons 10,5 et 11,5.
\end{enumerate}
On économise une antenne par rapport à \code{strategie\_1}. L'antenne en 12,5 est au-delà de la dernière maison : c'est permis, l'axe continue.

\Q{11}{2 points}{Écrire une fonction \code{strategie\_2}, sur le modèle de \code{strategie\_1}, qui applique cette nouvelle stratégie.}
\begin{reponse}
\begin{lstlisting}
def strategie_2(maisons, rayon):
    antennes = [Antenne(maisons[0].get_pos_maison() + rayon, rayon)]
    for m in maisons[1:]:
        if not antennes[-1].couvre(m):
            antennes.append(Antenne(m.get_pos_maison() + rayon, rayon))
    return antennes
\end{lstlisting}
\end{reponse}
\methode On recopie \code{strategie\_1} et on change seulement l'endroit où l'on pose une antenne : \code{position + rayon} au lieu de \code{position}, aux lignes 2 et 5.
\begin{itemize}
\item Ligne 2 : la première maison n'est jamais couverte au départ : on pose une antenne en \code{maisons[0]} $+$ \code{rayon}.
\item Lignes 3--4 : pour chaque maison suivante, on regarde si la \textbf{dernière} antenne posée, \code{antennes[-1]}, la couvre. Les maisons étant triées, les antennes précédentes sont plus à gauche et ne peuvent pas mieux faire.
\item Ligne 5 : sinon, nouvelle antenne en position de la maison $+$ \code{rayon}.
\end{itemize}
\textbf{Exemple.} Avec la rue de la question 9 et le rayon 2, \code{[a.get\_pos\_antenne() for a in strategie\_2(maisons, 2)]} vaut \code{[2, 7, 12.5]}. Le code a aussi été testé sur 2\,000 rues tirées au hasard : toutes les maisons sont couvertes, et il n'y a jamais plus d'antennes qu'avec \code{strategie\_1}.

\textbf{Pour aller plus loin.} Cet algorithme est \textbf{glouton} : à chaque étape, il fait le choix qui semble le meilleur sur le moment (aller le plus loin possible), sans revenir en arrière. Pour ce problème, on peut démontrer que ce choix donne le nombre minimal d'antennes.
\begin{erreur}
Écrire \code{Antenne(m + rayon, rayon)} : \code{m} est un objet \code{Maison}, pas un nombre, et Python s'arrête sur une \code{TypeError} (on ne peut pas additionner un objet et un nombre). Il faut sa position : \code{m.get\_pos\_maison()}.
\end{erreur}

\end{document}
